% !PRINTED 632
\weberpage{632}
%<<CONT
Stelle von $\omega'$ die natürliche Primzahl $p$ tritt, die durch $\mathfrak p$ theilbar
 ist, und wir reden danach von den \gesperrt{mit einer Zahl $\omega$
verwandten oder zu ihr fremden Primzahlen}.

\gesperrt{Eine Zahl $\omega$, die gar keine verwandten Primzahlen hat, ist
eine Einheit.}

Im Körper $\Omega_m$ (wie überhaupt in jedem Normalkörper) sind
alle conjugirten Zahlen $\omega_n$ mit denselben Primzahlen verwandt.

\begin{center}
\S.\ 174.

\gesperrt{Das Kummer'sche Theorem.}
\end{center}

Wir haben schon in der Kreistheilungstheorie (Bd.\,I, \S.\,169)
gewisse Zerlegungen der natürlichen Primzahlen in Factoren, die
aus Einheitswurzeln zusammengesetzt waren, kennen gelernt,
und es ist nun von Wichtigkeit, zu untersuchen, wie sich diese
Factoren zu den Primfactoren im Körper $\Omega_m$ verhalten.

Wir betrachten eine Primzahl $p$, die zum Exponenten 1
gehört, so dass also $m$ ein Theiler von $p-1$ ist, und setzen
\[
p - 1 = m m'.
\]

Aus den $p^{\mathrm{ten}}$ Einheitswurzeln $\varrho$ lassen sich nach Bd.\,I, \S.\,167
Perioden von je $m'$ Gliedern bilden, die, wenn $g$ eine primitive
Wurzel von $p$ (nicht, wie oben, von $m$) ist, den Ausdruck haben:
\weq{1}{
\begin{aligned}
\eta &= \varrho \phantom{^g} + \varrho^{g^m} + \varrho^{g^{2m}} + \cdots + \varrho^{g^{(m'-1)m}} \\
\eta_1 &= \varrho^g + \varrho^{g^{m+1}} + \varrho^{g^{2m+1}} + \cdots + \varrho^{g^{(m'-1)m+1}} \\
&\ \, \cdots\cdots\cdots\cdots\cdots\cdots\cdots\cdots\cdots\cdots \\
\eta_{m-1} &= \varrho^{g^{m-1}} + \varrho^{g^{2m-1}} + \varrho^{g^{3m-1}} + \cdots + \varrho^{g^{p-2}}.
\end{aligned}
}

Bezeichnen wir, wie im \S.\,169 des ersten Bandes, mit $\varepsilon$ eine
primitive $(p-1)^{\mathrm{te}}$ Einheitswurzel, so dass
\weq{2}{\varepsilon^{m'} = r,}
wie bisher, eine primitive $m^{\mathrm{te}}$ Einheitswurzel ist, so sind die
Resolventen
\begin{equation*}\tag*{3)}
(\varepsilon^\lambda, \varrho) = \varrho + \varepsilon^\lambda \varrho^g + \varepsilon^{2\lambda}\varrho^{g^2} + \cdots + \varepsilon^{(p-2)\lambda}\varrho^{g^{p-2}},
\end{equation*}
% [sic] printed "3)" without an opening parenthesis, unlike the other display numbers on this page
worin $\lambda$ jede ganze rationale Zahl sein kann, und speciell für
$\lambda = m'$:
\weq{4}{(r, \eta) = \eta + r\eta_1 + r^2\eta_2 + \cdots + r^{m-1}\eta_{m-1}.}
